\section{Proofs for closed elements and topological completeness}
\label{annex:p08-mt-topological-completeness}

\begin{annexitem}{annex:p08-mt.1}{thm:p08-mt-dual}{The direction of the residual}
\begin{proof}
Reverse the order of a Heyting algebra. For arbitrary \(a,b,c\),
\[
 (b\Rightarrow_H a)\leq_A c
 \ \Longleftrightarrow\ c\leq_H b\Rightarrow_H a
 \ \Longleftrightarrow\ c\wedge_H b\leq_H a
 \ \Longleftrightarrow\ a\leq_A b+_A c.
\]
This proves the Brouwerian residual law and the displayed operation.
The lattice bounds and finite operations reverse as stated. Reversing
the order a second time recovers the original lattice and residual,
because a residual is unique: it is the greatest admissible element in
the Heyting order, or the least in the Brouwerian order. Induction on the
formula grammar now gives its value in the same underlying set with
the translated operations. The designated bounds are exchanged, so
Heyting truth \(1_H\) is Brouwerian zero \(0_A\).
\end{proof}
\end{annexitem}

\begin{annexitem}{annex:p08-mt.2}{thm:p08-mt-closed}{Closed elements and the open complement}
\begin{proof}
The closure axioms imply monotonicity: if \(x\leq y\), then
\(y=x+y\), whence \(Cy=Cx+Cy\) and \(Cx\leq Cy\).
Closed elements include \(0,1\), are closed under finite joins, and are
closed under meets: \(C(a\cdot b)\leq Ca\cdot Cb=a\cdot b\), while
extensivity gives the opposite inequality. Thus they inherit a bounded
distributive lattice. Their displayed difference is closed by idempotence.
For closed \(c\), monotonicity and extensivity give
\[
 C(a\cdot(-b))\leq c
 \ \Longleftrightarrow\ a\cdot(-b)\leq c
 \ \Longleftrightarrow\ a\leq b+c.
\]
The last equivalence is Boolean algebra, and proves residuation in both
directions. This is precisely the argument of the source's Theorem~1.14.

For subsets, Boolean difference is set difference and \(C\) is ordinary
closure. Complementation interchanges closed sets and open sets,
reverses inclusion, and exchanges finite union and intersection. Let
\(U=X\setminus a\), \(V=X\setminus b\). The Heyting implication has
the reversed difference order, so its open representative is
\[
 X\setminus\operatorname{cl}(b\setminus a)
 =\operatorname{int}((X\setminus b)\cup a)
 =\operatorname{int}(V\cup(X\setminus U)).
\]
This also proves directly that the modern open implication is not the
unreversed closed difference.
\end{proof}
\end{annexitem}

\begin{annexitem}{annex:p08-mt.3}{thm:p08-mt-extension}{Normal forms, well-defined closure, and uniqueness}
\begin{proof}
Boolean disjunctive normal form writes every element of the subalgebra
generated by \(e[L]\) as a finite join of terms \(e(a)\cdot(-e(b))\).
Indeed finite positive factors combine to one \(e(a)\), finite negative
factors to \(-e(b)\); empty positive and negative lists use \(1\) and
\(0\). Consequently the prescription is defined on all of \(K\).

For a chosen normal form \(x=\sum_i e(a_i)\cdot(-e(b_i))\), put
\(d=\sum_i(a_i\mathbin{\dot{-}}b_i)\). For each \(c\in L\),
\[
 \begin{split}
 x\leq e(c)
 &\Longleftrightarrow
  e(a_i)\cdot(-e(b_i))\leq e(c)\quad\hbox{for every }i\\
 &\Longleftrightarrow a_i\leq b_i+c\quad\hbox{for every }i\\
 &\Longleftrightarrow a_i\mathbin{\dot{-}}b_i\leq c
     \quad\hbox{for every }i\\
 &\Longleftrightarrow d\leq c.
 \end{split}
\]
Also \(x\leq e(d)\), since
\(a_i\leq b_i+(a_i\mathbin{\dot{-}}b_i)\). Thus \(e(d)\) is the
least element of \(e[L]\) above \(x\). Least elements are unique, so a
second normal form gives exactly the same \(Cx=e(d)\); this resolves
the construction's essential well-definedness obligation.

The empty normal form gives \(C0=0\). Extensivity was just proved.
Every \(e(a)\) is its own least majorant, so \(Ce(a)=e(a)\), and
\(CCx=Cx\). Concatenating normal forms gives
\(C(x+y)=Cx+Cy\). Therefore \((K,C)\) is a closure algebra. Its
closed elements are exactly \(e[L]\): each is fixed, and every fixed
point equals a value of \(C\) in \(e[L]\). The one-term normal form
gives \(C(e(a)\cdot(-e(b)))=e(a\mathbin{\dot{-}}b)\).

For uniqueness, let \(u:L\to L'\) be a Brouwerian isomorphism and
let \(e'\) be another such embedding. A signed Boolean monomial
\(\prod_i e(a_i)\cdot\prod_j(-e(b_j))\) is zero exactly when
\(\prod_i a_i\leq\sum_j b_j\). These lattice inequalities are
preserved and reflected by \(u\), including empty lists. By normal
form, every Boolean term is zero in one extension exactly when its
image is zero in the other. Equality of two terms reduces to their
Boolean symmetric difference being zero. Hence sending \(e(a)\) to
\(e'(u(a))\) is a well-defined Boolean isomorphism of the generated
algebras. It commutes with closure by the least-majorant property (or
the displayed normal-form formula), and is unique because its values
on the Boolean generators determine all other values.

For the assertion about a starting embedding, apply the already proved
prime-filter representation, Theorem~\ref{thm:p10-heyting-prime-representation},
to \(H=L^{\mathrm{op}}\). Under the stated Boolean prime ideal theorem,
its injective map \(\beta_H\) into upsets gives the bounded lattice
embedding \(e(a)=W\setminus\beta_H(a)\) into \(\mathcal P(W)\).
This is a modern way to supply the Boolean embedding assumed by the
source's Theorem~1.15, not an attribution of that prime-filter proof to
1946. The finite construction in item~\ref{annex:p08-mt.6} removes this
choice dependency in the completeness argument below.
\end{proof}
\end{annexitem}

\begin{annexitem}{annex:p08-mt.4}{thm:p08-mt-finite}{Finite lattice, residual, and subterm preservation}
\begin{proof}
Every bounded lattice term in the \(n\) elements of \(M\) is, by
distributivity, a join of meets of subsets of \(M\). Empty meets and
joins supply \(1\) and \(0\). There are at most \(2^n\) such meets
and at most \(2^{2^n}\) joins of them. This proves finiteness and the
bound, including \(n=0\).

For \(a,b\in D\), the finite nonempty set
\(S=\{z\in D:a\leq b+z\}\) contains \(1\). Put
\(a\mathbin{\dot{-}}_D b=\prod S\). Distributivity gives
\(b+\prod S=\prod_{z\in S}(b+z)\), so \(a\leq b+\prod S\).
Thus the meet belongs to \(S\) and is its least element; this establishes
the residual law in \(D\). If the original difference
\(d=a\mathbin{\dot{-}}_A b\) lies in \(D\), it belongs to \(S\)
and is below every member of \(S\), by original residuation. Hence the
new difference equals \(d\).

Include in \(M\) the values of every subterm of each listed term under
the given assignment. Interpret its variables by their original values.
Induction along the finite term syntax preserves joins, meets, and
bounds because these are inherited in \(D\). At a difference node,
the values of both inputs and of the output are in \(M\subseteq D\),
so the preceding agreement clause preserves that node too. This proves
all promised term values, but does not assert agreement on unlisted
differences. Taking a nonzero term value in \(M\) preserves that
nonzeroness, because \(D\) uses the original zero.
\end{proof}
\end{annexitem}

\begin{annexitem}{annex:p08-mt.5}{thm:p08-mt-minimal}{Relativization and the minimal counteralgebra}
\begin{proof}
For \(x,y,z\geq a\),
\[
 a+(x\mathbin{\dot{-}}y)\leq z
 \ \Longleftrightarrow\ x\mathbin{\dot{-}}y\leq z
 \ \Longleftrightarrow\ x\leq y+z.
\]
So \(A_a\) is a Brouwerian algebra with zero \(a\). The map
\(\rho_a(x)=a+x\) preserves join, meet by distributivity, and both
bounds. For every \(z\geq a\),
\[
 \rho_a(x)\leq\rho_a(y)+z
 \ \Longleftrightarrow\ x\leq y+z
 \ \Longleftrightarrow\ a+(x\mathbin{\dot{-}}y)\leq z.
\]
Uniqueness of the least admissible element in \(A_a\) therefore proves
\(\rho_a(x\mathbin{\dot{-}}y)
 =\rho_a(x)\mathbin{\dot{-}}_a\rho_a(y)\).

Finite preservation provides at least one finite algebra and assignment
where the chosen term is nonzero. Choose one of smallest cardinality;
this is a least natural number, not a selection over an infinite family.
If it is not well-connected, choose nonzero \(a,b\) with
\(a\cdot b=0\). Both relativized algebras have smaller cardinality,
since their carriers omit the original zero. The map
\(x\mapsto(a+x,b+x)\) is injective: distributivity gives
\((a+x)\cdot(b+x)=x+(a\cdot b)=x\).
If a nonzero term value \(d\) became zero in both factors, then
\(a+d=a\) and \(b+d=b\), implying \(d\leq a\cdot b=0\), a
contradiction. One factor therefore has the same term nonzero under
the image assignment; term values commute with its homomorphism by
syntax induction. This contradicts minimality and proves well-connectedness.
The dual formula claim follows from item~\ref{annex:p08-mt.1}.
\end{proof}
\end{annexitem}

\begin{annexitem}{annex:p08-mt.6}{thm:p08-mt-finite-space}{Finite representation without an unproved duality step}
\begin{proof}
In a finite distributive lattice, every nonzero element is a join of
nonzero join-irreducibles below it: if it is not itself irreducible,
split it into a join of strictly smaller elements and use induction in
the finite order. A join-irreducible \(j\) is join-prime, because
\(j\leq a+b\) implies
\(j=(j\cdot a)+(j\cdot b)\), so \(j\leq a\) or \(j\leq b\).
It follows that \(e\) preserves joins, meets, and bounds. If
\(a\not\leq b\), the decomposition of \(a\) supplies some
\(j\leq a\) not below \(b\); thus \(e\) reflects order and is
injective, with the promised separation witness.

Every downset \(S\subseteq J(L)\) equals \(e(\sum S)\). One inclusion
is immediate; for the other, join-primality puts an irreducible below
one of the finitely many members of \(S\), and downward closure puts
it in \(S\). This includes the empty set. Hence \(e[L]\) is the
entire closed-set lattice of the upper Alexandrov topology.
The least downset containing \(e(a)\setminus e(b)\) is
\(\downarrow(e(a)\setminus e(b))\). For a downset \(e(c)\), it is
contained in \(e(c)\) exactly when
\(e(a)\setminus e(b)\subseteq e(c)\), exactly when \(a\leq b+c\).
By residuation and uniqueness of this least set it equals
\(e(a\mathbin{\dot{-}}b)\). This proves difference preservation
and the identification of the normal-form closure.

For each \(j\in J(L)\),
\[
 \{j\}=e(j)\setminus\bigcup_{k<j}e(k).
\]
Thus the Boolean-generated algebra contains every singleton, and by
finiteness is all of \(\mathcal P(J(L))\). This explicitly supplies
the Boolean extension used in the finite case.

If the lattice is nontrivial, \(J(L)\) is nonempty. If it had two
different minimal points, their principal downsets would be disjoint
nonempty closed sets, contradicting well-connectedness. A finite poset
with a unique minimal point has that point below every element: descend
below any chosen element to a minimal one. Hence that point is least.
If the lattice is trivial, \(J(L)=\varnothing\); the topology and its
closed-set lattice are also trivial. Such an algebra cannot witness
a nonzero term.
\end{proof}
\end{annexitem}

\begin{annexitem}{annex:p08-mt.7}{thm:p08-mt-formulas}{Formula-by-formula preservation}
\begin{proof}
Atoms satisfy the equation by definition, and bottom does because
\(e(1)=J(L)\). Complementation exchanges union and intersection, so
the clauses for conjunction and disjunction follow from their term
translations and the preservation of the lattice operations. For implication,
put \(a=t_B\), \(b=t_D\). Its translated closed value is
\(b\mathbin{\dot{-}}a\), and item~\ref{annex:p08-mt.6} gives
\[
 \begin{split}
 J(L)\setminus e(b\mathbin{\dot{-}}a)
 &=J(L)\setminus\operatorname{cl}(e(b)\setminus e(a))\\
 &=\operatorname{int}((J(L)\setminus e(b))\cup e(a))\\
 &=\operatorname{int}(\nu(D)\cup(J(L)\setminus\nu(B))).
 \end{split}
\]
This is exactly the prescribed topological implication and completes
the syntax induction. Because \(e\) is injective and \(e(0)=\varnothing\),
the designated-value equivalence and the pointwise separating witness
follow in both directions.

In the upper Alexandrov topology \(\uparrow w\) is the least open
neighborhood of \(w\). Thus \(w\) belongs to the interior of a set
\(S\) iff \(\uparrow w\subseteq S\). Applying this to
\(S=(W\setminus\nu(B))\cup\nu(D)\) proves the stated future-world
clause, including reflexivity at \(w\).
\end{proof}
\end{annexitem}

\begin{annexitem}{annex:p08-mt.8}{thm:p08-mt-cantor}{The local coloring property}
\begin{proof}
At stage \(m\) of the Cantor construction, \(K\) is covered by
\(2^m\) closed basic intervals of length \(3^{-m}\). For \(x\in K\)
and an open neighborhood \(O\) of \(x\), choose \(\epsilon>0\) with
\((x-\epsilon,x+\epsilon)\subseteq O\), and choose \(m\) with
\(3^{-m}<\epsilon\). Any basic interval of stage \(m\) containing
\(x\) is wholly inside \(O\): every point of that interval is within
\(3^{-m}\) of \(x\). For every \(n>m\), a descendant interval of
stage \(n-1\) has a middle third removed at generation \(n\), and
this entire open removed interval is inside the chosen stage-\(m\)
interval. The integers \(n>m\) run through every residue class modulo
\(k\). Choosing one such \(n\) for each color proves the claim.
Nothing in the argument requires \(x\) to lie inside a basic interval
rather than at its endpoint; it therefore includes \(0\) and \(1\).
\end{proof}
\end{annexitem}

\begin{annexitem}{annex:p08-mt.9}{thm:p08-mt-realization}{Constructing the continuous open surjection}
\begin{proof}
Induct on the cardinality of \(W\). For a singleton, the constant map
is continuous, open, and onto: the image of a nonempty open set is the
unique nonempty open set of the singleton.

Otherwise let \(s_1,\ldots,s_k\) be all immediate successors of \(r\).
Every \(q>r\) lies above some \(s_i\), by taking a minimal element
of the nonempty finite set \(\{u:r<u\leq q\}\).
Each \(W_i=\uparrow s_i\) is an upper subspace with least point
\(s_i\), and has smaller cardinality. By induction choose a continuous
open surjection \(f_i:\mathbb R\to W_i\).
For each bounded complementary interval \(I\) of color \(i\), choose
an increasing homeomorphism \(h_I:I\to\mathbb R\), for example
\(h_I(x)=\log((x-a)/(b-x))\) when \(I=(a,b)\).
For the exterior intervals one can use
\(h_{(-\infty,0)}(x)=-\log(-x)\) and
\(h_{(1,\infty)}(x)=\log(x-1)\), respectively. On a bounded interval
of color \(i\), define \(f=f_i\circ h_I\); on the two exterior
intervals use \(f_1\circ h_I\); on \(K\) define \(f(x)=r\).
These disjoint pieces partition \(\mathbb R\), so \(f\) is well
defined. Every \(q>r\) is reached in an interval of some color \(i\)
with \(s_i\leq q\), and \(r\) is reached on \(K\); thus \(f\) is
surjective.

For continuity, take an upset \(U\subseteq W\). If it contains
\(r\), it equals \(W\), so its inverse image is \(\mathbb R\).
Otherwise its inverse image is a union of subsets of complementary
intervals. On each interval, this subset is open by continuity of
\(f_i\circ h_I\), since \(U\cap W_i\) is open in \(W_i\).
An interval is itself open in \(\mathbb R\), so the union is open.

For openness, let \(O\subseteq\mathbb R\) be open. If \(O\cap K\)
is empty, then \(f[O]\) is a union of images of open subsets of the
complementary intervals. Each image is an upset in \(W_i\) by induction,
and hence in \(W\), because \(W_i\) itself is an upset. Thus the
image is open. If \(x\in O\cap K\), the coloring lemma supplies,
for each \(i\), a whole interval \(I_i\subseteq O\) of color \(i\).
Its image is all of \(W_i\); together these images cover every point
above \(r\), while \(f(x)=r\). Hence \(f[O]=W\), again open.
This proves openness at the root as well as on all interval pieces,
and completes the induction.
\end{proof}
\end{annexitem}

\begin{annexitem}{annex:p08-mt.10}{thm:p08-mt-pullback}{Interior, closure, and all operations}
\begin{proof}
Continuity makes \(f^{-1}[\operatorname{int}_Y S]\) an open subset
of \(f^{-1}[S]\), yielding one inclusion. Conversely, if an open
\(O\subseteq X\) lies in \(f^{-1}[S]\), then \(f[O]\) is open by
openness and lies in \(S\). Thus
\(O\subseteq f^{-1}[\operatorname{int}_Y S]\). Taking the union of
all such \(O\) proves the interior equality. Applying it to
\(Y\setminus S\), and using the closure--interior complement identity
in both spaces, proves the closure equality.

Inverse images preserve the empty set, the entire carrier, intersections,
unions, and complements. Syntax induction consequently preserves
atomic valuations, bottom, conjunction, and disjunction. At implication,
\[
 \begin{split}
 f^{-1}[\operatorname{int}_Y((Y\setminus\nu(B))\cup\nu(D))]
 &=\operatorname{int}_X((X\setminus f^{-1}[\nu(B)])
                      \cup f^{-1}[\nu(D)]).
 \end{split}
\]
This is the implication of the pulled-back valuation. For closed
terms, the corresponding difference clause is
\(f^{-1}[\operatorname{cl}_Y(A\setminus B)]
 =\operatorname{cl}_X(f^{-1}[A]\setminus f^{-1}[B])\).
Induction gives every Brouwerian term as well.
If \(f\) is surjective, \(f^{-1}[S]=X\) iff \(S=Y\), and
\(f^{-1}[S]=\varnothing\) iff \(S=\varnothing\). In particular
each failure point \(y\notin\nu(B)\) has an inverse-image point
\(x\notin f^{-1}[\nu(B)]\). These prove preservation and reflection
with their distinct hypotheses.
\end{proof}
\end{annexitem}

\begin{annexitem}{annex:p08-mt.11}{thm:p08-mt-topological-completeness}{Soundness and the constructed countervaluation}
\begin{proof}
For every space, its opens form a Heyting algebra: the residual follows
from item~\ref{annex:p08-mt.2}, or directly from
\(Z\cap U\subseteq V\) iff
\(Z\subseteq\operatorname{int}((X\setminus U)\cup V)\) for open
\(Z\). The soundness direction of
Theorem~\ref{thm:p10-heyting-completeness} therefore proves that IPL
theorems are valid on every space, including the finite spaces and
each positive-dimensional Euclidean space. The empty space causes no
exception to soundness: its open algebra is trivial and every formula
has value its entire (empty) carrier.

Now suppose \(\nvdash_{\mathrm{IPL}}B\). The same algebraic theorem's
Lindenbaum construction gives a Heyting algebra and an assignment with
\(v(B)\ne1\). In its Brouwerian order dual, the translated term
\(t_B\) is nonzero by item~\ref{annex:p08-mt.1}.
Item~\ref{annex:p08-mt.4} preserves all its subterm values in a finite
algebra, and item~\ref{annex:p08-mt.5} supplies a finite well-connected
algebra \(A\) with an assignment still making \(t_B\ne0\).
It is nontrivial, so item~\ref{annex:p08-mt.6} supplies a nonempty
finite space \(W=J(A)\) with least point \(r\).
The open valuation of item~\ref{annex:p08-mt.7} gives
\(\nu(B)=W\setminus e(t_B)\ne W\). More specifically,
\(e(t_B)\) is a nonempty downset, so it contains the least point
\(r\); thus \(r\notin\nu(B)\). This proves the failure on a finite
upper Alexandrov space and provides an explicit separating point.

Pull this valuation back along the map
\(f:\mathbb R\to W\) of item~\ref{annex:p08-mt.9}. Every formula
value is preserved by item~\ref{annex:p08-mt.10}. For every \(x\in K\),
\(f(x)=r\notin\nu(B)\), so \(x\) is a failure point. The pulled-back
valuation therefore witnesses failure on \(\mathbb R\).
For each fixed integer \(d\geq1\), the first-coordinate projection
\(\pi:\mathbb R^d\to\mathbb R\) is continuous, open, and onto.
Openness follows because a basic open product has an open first-coordinate
image and every open set is a union of such products; continuity and
surjectivity are immediate from the product topology. Pullback along
\(\pi\) preserves the formula again, with every point over \(K\) a
failure point. Thus nonderivability fails each of (ii)--(v), proving
all the equivalences.

All formulas and subterm lists are finite. The Lindenbaum quotient and
the explicitly finite lattice constructions require no prime-filter
selection. The real maps use only finite recursion and the specified
Cantor intervals. Hence no Boolean prime ideal theorem is hidden in
this countervaluation proof. None of these arguments constructs a
valuation satisfying an arbitrary infinite premise theory on a fixed
Euclidean space, so no such strong-completeness assertion follows.
\end{proof}
\end{annexitem}

\begin{annexitem}{annex:p08-mt.12}{thm:p08-mt-functional-free}{Separating any failed Brouwerian identity in Euclidean space}
\begin{proof}
Any identity valid in all Brouwerian algebras is valid in the Euclidean
closed-set algebra, because that is a Brouwerian algebra by
item~\ref{annex:p08-mt.2}. For the converse, suppose \(s=t\) fails
under an assignment in an algebra \(A\). By residuation,
\(a\mathbin{\dot{-}}b=0\) iff \(a\leq b\). Therefore the term
\[
 q=(s\mathbin{\dot{-}}t)+(t\mathbin{\dot{-}}s)
\]
has a nonzero value: otherwise both differences are zero, giving both
orders and hence equality of the two values.
Apply finite preservation and the minimal well-connected reduction to
\(q\), including all its subterms. This produces a finite rooted
space and a closed assignment where \(q\) is a nonempty downset,
using item~\ref{annex:p08-mt.6}. Pull it back along the constructed
continuous open surjection from \(\mathbb R\). Closedness and all
term operations are preserved by item~\ref{annex:p08-mt.10}, and
nonemptiness is reflected by surjectivity. Thus \(q\ne\varnothing\)
on \(\mathbb R\), so \(s\ne t\) there. Projection pullback gives
the same separating closed assignment in \(\mathbb R^d\) for every
\(d\geq1\). This proves exactly identity-level functional freeness.
It does not select generators or prove freeness of any closure algebra;
those are different algebraic properties.
\end{proof}
\end{annexitem}

\begin{annexitem}{annex:p08-mt.13}{thm:p08-mt-boundaries}{Examples, excluded cases, and the necessity of map conditions}
\begin{proof}
In the two-point upper topology, the closed complement of \(\{s\}\)
is \(\{r\}\), whose interior is empty. Consequently
\(\nu(\neg p)=\operatorname{int}(\{r\})=\varnothing\), and the
excluded-middle value is \(\{s\}\), missing \(r\).
The closed algebra has carrier \(\varnothing,\{r\},W\). For its
middle element \(a=\{r\}\),
\(1\mathbin{\dot{-}}a=\operatorname{cl}(W\setminus\{r\})
=\downarrow\{s\}=W\), so
\(a\cdot(1\mathbin{\dot{-}}a)=a\ne\varnothing\).
The ordinary complement \(\{s\}\) is not closed, so substituting it
for the Brouwerian complement would leave the closed carrier.
On \(\mathbb R\), the complement of the dense open set
\(U=\mathbb R\setminus\{0\}\) is \(\{0\}\), also with empty
interior. Hence \(\neg U=\varnothing\) and
\(\neg\neg U=\operatorname{int}(\mathbb R)=\mathbb R\), whereas
\(U\cup\neg U=U\). These are concrete non-Boolean open algebras.

On a discrete space every subset is open and interior is the identity.
Thus open implication is \((X\setminus U)\cup V\), negation is
ordinary complement, and excluded middle is \(X\). The singleton
\(\mathbb R^0\) is discrete, so validates this formula that fails in
the preceding finite and real spaces. It cannot characterize IPL.

For the continuous but nonopen counterexample, let \(X\) be the
discrete two-point set and \(Y\) its upper Sierpi\'nski topology, and
let \(f:X\to Y\) be the identity. Every inverse image of a target
open is open in \(X\), so \(f\) is continuous and onto. The image
of the source open \(\{r\}\) is not open in \(Y\), so it is not
open. For \(S=\{r\}\),
\(f^{-1}[\operatorname{int}_Y S]=\varnothing\), whereas
\(\operatorname{int}_X(f^{-1}[S])=\{r\}\). This explicitly falsifies
the implication-preserving interior equality without openness.
In the reverse direction, the identity from \(Y\) to the discrete
space is open, but not continuous: the inverse image of the discrete
open \(\{r\}\) is not open in \(Y\).
Finally, the inclusion of the open subspace \(\{s\}\) in \(Y\) is
continuous and open but not onto; it misses the only failure point of
excluded middle in the first valuation. Its pulled-back valuation
validates that formula on the singleton. This proves that surjectivity
is needed for reflection, not merely for convenience in the construction.
\end{proof}
\end{annexitem}
